% Discussion v2 (post v2 round-1 review fixes).

\paragraph{Headline.}
On the unified ESMO+ASCO+NCCN HR+/HER2- mBC decision tree (25 nodes) against a 66-edge trial-DAG, the strict-tolerance evidence-free decision-point ratio was 0.48 (Clopper-Pearson 95\% CI 0.28--0.69), and the pre-registered one-sided exact binomial test of $H_0: \mathrm{EFDPR} \le 0.25$ rejected the null at $\alpha = 0.05$ with $P = 0.011$. ESCAT-aligned and liberal tolerance gave identical estimates and rejections ($P = 0.011$). The pilot-to-production scaling from 16 nodes (pilot, $P = 0.37$) to 25 nodes (production, $P = 0.011$) increased statistical power and surfaced NCCN-specific TROP2-ADC nodes, ASCO-specific special-population nodes, and post-CDK4/6i biomarker-stratified positions as the largest evidence gaps. Subset sensitivity shows that the rejection is driven primarily by the NCCN-citing nodes ($P = 0.020$); ESMO-only ($P = 0.19$) and ASCO-only ($P = 0.31$) subsets do not individually reject the null in this corpus, a transparency point we emphasise in the present discussion.

\paragraph{What the framework actually measures.}
The strict concordance algorithm is a deliberately conservative model: it counts as evidence-free any guideline node where no single pivotal trial enrolled exactly the (state, biomarker) population the guideline assumes. This is by design. The point is to make explicit, on one diagram, where the guideline tree is supported by direct trial chains versus where it is supported by composition across non-overlapping trials. Composition-across-non-overlap is a routine and often unavoidable practice for rapidly evolving fields; it is not illegitimate, but it has not been counted before.

\paragraph{Where the gap is largest.}
Three structural patterns dominate the 12 evidence-free nodes under strict concordance: (i) NCCN-specific TROP2-ADC nodes (G20 sacituzumab govitecan for post-endo+post-CDK4/6i HR+/HER2-, G21 datopotamab deruxtecan for post-chemo HER2-low, G23 sacituzumab govitecan for late-line HR+/HER2-) are recent and inherit the maturation gap between trial readout year and guideline adoption; (ii) ASCO-specific special-population nodes (G17 visceral-crisis 1L chemotherapy, G22 indolent-disease single-agent endocrine, G24 post-menopausal indolent monotherapy, G25 post-everolimus salvage) cover patient populations rarely enrolled in modern phase 3 trials; (iii) post-CDK4/6i biomarker-stratified positions (G9 post-CDK4/6i no-actionable, G13 post-CDK4/6i+post-endo everolimus) remain evidence-free because BOLERO-2 enrolled a post-AI rather than post-CDK4/6i population. After adding BYLieve (post-CDK4/6i PIK3CAmut, NCT03056755) and TROPION-Breast01 HER2-low subgroup readout, the post-CDK4/6i ESR1mut (G5) and PIK3CAmut (G7) nodes are now supported under strict tolerance --- a substantive change from the v2 round-0 draft that flagged G7 as evidence-free.

\paragraph{Operationalization discordance.}
On 82 trials and 3,240 trial pairs the prior-CDK4/6i inclusion variable had ODI 0.64 (bootstrap 95\% CI 0.62--0.66) --- the most consequential single-variable heterogeneity. Trials variously require, exclude, or permit prior CDK4/6i, and post-CDK4/6i subgroups derived from permit-prior-CDK4/6i populations cannot be cleanly combined with post-CDK4/6i subgroups from require-prior-CDK4/6i populations. ESR1 mutation definitions (ODI 0.38) showed the second-largest discordance, driven by tissue vs plasma sampling and the variant call set. These operational findings affect not only this framework's concordance algorithm but also any network-meta-analysis that attempts to pool across post-CDK4/6i trials.

\paragraph{Three caveats.}
First, the unit of inference is the guideline decision node, not the trial. With $n = 25$ nodes the test is moderately powered ($\sim 42\%$ at $p_1 = 0.40$, $\sim 79\%$ at $p_1 = 0.50$). The pre-registered exact-binomial test rejected at $P = 0.011$, but the rejection is driven by the NCCN-citing subset and a fully-powered confirmatory study would pool decision nodes across multiple solid tumours. Second, the LLM-extraction pipeline (Annotator A: Claude; Annotator B: Codex/GPT-5) achieved high inter-rater agreement post-adjudication (mean PABAK 0.99) but the pre-registered Cohen's $\kappa \ge 0.70$ gate failed pre-adjudication on \texttt{post\_endo} ($\kappa = 0.40$) and on \texttt{akt\_path} ($\kappa = 0.00$ due to the kappa paradox); we report both Cohen's $\kappa$ and PABAK, and disclose the failure of the Cohen's $\kappa$ gate for the paradox case. The v2 round-1 internal review surfaced two pivotal trials missing from the systematic corpus (BYLieve, SERENA-6 with the previously confused NCT04711252 vs NCT04964934 identifier) and two extraction errors (CAPItello-291 and DESTINY-Breast06 post-CDK4/6i subgroup-readout flags, TROPION-Breast01 HER2-low subgroup-readout flag); after correction the strict-tolerance EFDPR moved from 0.60 to 0.48 and the primary test continued to reject at $P = 0.011$ rather than $P = 0.0002$. Third, the by-guideline-source sensitivity shows that ESMO alone and ASCO alone do not individually reject the null; the primary rejection therefore depends on the unified-tree denominator, which is a defensible but consequential choice that the manuscript and supplement disclose explicitly.

\paragraph{Three positives.}
First, the framework is fully reproducible from the public ClinicalTrials.gov API and the public guideline documents. The 874-candidate $\to$ 82-trial filter pipeline is committed; the dual-annotator extraction protocol and per-trial outputs are released; the adjudication rules (17 total) and post-adjudication kappas are published. Second, the framework is tumour-agnostic: the (state, biomarker) node space, the concordance algorithm, and the pre-registered binomial test transfer to any solid tumour whose guideline tree composes across non-overlapping pivotal trials; mBC is the first application. Third, the pre-registered design with two-annotator validation against an explicit protocol, plus an internal multi-round adversarial review that surfaced and corrected concrete extraction and corpus errors before release, provides a template for LLM-extraction studies that aspire to confirmatory inference.

\paragraph{Limitations.}
The 82-trial corpus is a snapshot at freeze date 2026-05-16 and excludes trials primarily completed after 2026. The guideline encoding fixes the ESMO 2024 update, ASCO 2022 living guideline, and NCCN 2024 update; ASCO and NCCN both update more frequently than annual and the EFDPR overlay should be refreshed at each update. Eligibility coding by LLM remains imperfect for highly inferential fields (\texttt{post\_endo} pre-adjudication $\kappa = 0.40$); production use should retain the dual-annotator + protocol-adjudication discipline. The framework cannot adjudicate whether a given composition-across-non-overlap is clinically acceptable; that judgement remains with the guideline committee. Finally, the v1.0.0 16-trial pilot (preserved at tag \texttt{v1.0.0}) gave a non-rejecting result ($P = 0.37$); the v2 production result rejects, but the change is consistent with statistical power from a larger $n$ and a fully-disclosed corpus expansion, not with re-analysis of the same data.

\paragraph{Outlook.}
The framework, schema, decision-tree encoding, and dual-annotator extraction protocol are released to enable extension by guideline committees, trial sponsors, and meta-research groups. A community-maintained Treatment-Sequencing Evidence Atlas, refreshed at every ESMO, ASCO, and NCCN HR+/HER2- mBC guideline release, could provide a continuously-updated EFDPR overlay across solid tumours. Immediate priorities are the NSCLC (EGFR-mutant) and mCRC (ctDNA-guided) applications, a fully-LLM-extracted census of HR+/HER2- mBC trials with inter-rater agreement at production scale, and an interactive web viewer that lets a clinician click any guideline decision node and see, in one screen, every trial-DAG edge that supports it and every edge that does not.
